\documentclass[11pt]{article}
\usepackage[margin=1in]{geometry}
\usepackage{amsmath,amssymb,amsthm,mathtools}
\usepackage[T1]{fontenc}
\usepackage[utf8]{inputenc}
\usepackage{enumitem}
\usepackage{booktabs,longtable,array}
\usepackage{hyperref}
\usepackage{xcolor}
\usepackage{microtype}
\hypersetup{colorlinks=true, linkcolor=blue!50!black, citecolor=blue!50!black, urlcolor=blue!50!black}

\newtheorem{theorem}{Theorem}[section]
\newtheorem{proposition}[theorem]{Proposition}
\newtheorem{lemma}[theorem]{Lemma}
\newtheorem{corollary}[theorem]{Corollary}
\newtheorem{claim}[theorem]{Claim}
\theoremstyle{definition}
\newtheorem{definition}[theorem]{Definition}
\newtheorem{assumption}[theorem]{Assumption}
\theoremstyle{remark}
\newtheorem{remark}[theorem]{Remark}

\newcommand{\KL}{\mathrm{KL}}
\newcommand{\Ent}{\mathrm{Ent}}
\newcommand{\Prob}{\mathsf{P}}
\newcommand{\E}{\mathbb{E}}
\newcommand{\R}{\mathbb{R}}
\newcommand{\N}{\mathbb{N}}
\newcommand{\Pcal}{\mathcal{P}}
\newcommand{\Mcal}{\mathcal{M}}
\newcommand{\Qcal}{\mathcal{Q}}
\newcommand{\Zcal}{\mathcal{Z}}
\newcommand{\Fcal}{\mathcal{F}}
\newcommand{\Gcal}{\mathcal{G}}
\newcommand{\W}{\mathcal{W}}
\newcommand{\one}{\mathbf{1}}
\newcommand{\dd}{\,d}
\newcommand{\supp}{\operatorname{supp}}
\newcommand{\argmin}{\operatorname*{arg\,min}}
\newcommand{\argmax}{\operatorname*{arg\,max}}
\newcommand{\esssup}{\operatorname*{ess\,sup}}
\newcommand{\ri}{\operatorname{ri}}
\newcommand{\dom}{\operatorname{dom}}
\newcommand{\Law}{\operatorname{Law}}
\newcommand{\Var}{\operatorname{Var}}
\newcommand{\Cov}{\operatorname{Cov}}
\newcommand{\diag}{\operatorname{diag}}
\newcommand{\Id}{\mathrm{Id}}
\newcommand{\Lip}{\operatorname{Lip}}
\newcommand{\osc}{\operatorname{osc}}
\newcommand{\diam}{\operatorname{diam}}
\newcommand{\proj}{\operatorname{proj}}
\newcommand{\e}{\mathrm{e}}

\newenvironment{argument}{\paragraph{Argument.}\begin{enumerate}[leftmargin=2em]}{\end{enumerate}}
\newenvironment{proofoutline}{\paragraph{Proof architecture.}\begin{enumerate}[leftmargin=2em]}{\end{enumerate}}


% Source-faithfulness apparatus used by the paper reconstructions.
\newcommand{\sourceanchor}[1]{\par\smallskip\noindent\textbf{Source anchor.} #1\par\smallskip}
\newcommand{\sourcecomment}[1]{\begin{remark}#1\end{remark}}
\newenvironment{sourceguide}
  {\begin{description}[style=nextline,leftmargin=3.2cm,labelwidth=3cm]}
  {\end{description}}

% Backward-compatible environments used by some of the older reconstructions.
\newenvironment{distilled}[1][]{\begin{theorem}[#1]}{\end{theorem}}
\newenvironment{distillation}{\begin{enumerate}[leftmargin=2em]}{\end{enumerate}}

\newenvironment{commentarynotes}{\begin{remark}\begin{enumerate}[leftmargin=2em]}{\end{enumerate}\end{remark}}

\title{Wang--Gao--Xie, \emph{Sinkhorn DRO}: Theorem 1 (Strong Duality), in full}
\author{Distilled literature note --- the baseline our Lean proof is compared against}
\date{Source: \texttt{papers/WangGaoXie2023-2109.11926.pdf} (arXiv v-of-record), pp.~6--8, 11--13;
supplement \texttt{papers/WangGaoXie2023-supplement-opre.2023.0294.sm2.pdf}}
\begin{document}
\maketitle

\section*{Purpose of this note}
This note records the mathematics \textbf{as it appears in \cite{WangGaoXie2025}}, not as we
formalized it.  It is the independent baseline against which
\texttt{WangGaoXie2023.strong\_duality} and \texttt{Drsb.sdrsb\_strong\_duality} are to be diffed.
Where our Lean development deviates --- and it does, in one identified place --- the deviation is
recorded in \S\ref{sec:diff}, with the paper's claim taken as authoritative unless we have refuted it.

\section{Setting and notation (paper's own)}

$\Zcal$ is a measurable space, $\widehat{\Prob}\in\Pcal(\Zcal)$ the nominal distribution,
$f:\Zcal\to\R\cup\{\infty\}$ the loss, $c:\Zcal\times\Zcal\to\R$ the transport cost,
$\epsilon>0$ the entropic regularization parameter, and $\mu,\nu\in\Mcal(\Zcal)$ two reference
measures.  The paper writes $\rho$ for the ambiguity radius.

\paragraph{Definition 1 (Sinkhorn distance).}
For $P,Q\in\Pcal(\Zcal)$ with $P\ll\mu$, $Q\ll\nu$ and $\epsilon\ge 0$,
\[
  \W_\epsilon(P,Q)\;=\;\inf_{\gamma\in\Gamma(P,Q)}
    \Big\{\, \E_{(x,y)\sim\gamma}[c(x,y)] \;+\; \epsilon\, H(\gamma\mid\mu\otimes\nu) \,\Big\},
  \qquad
  H(\gamma\mid\mu\otimes\nu)=\E_{(x,y)\sim\gamma}\Big[\log\tfrac{d\gamma(x,y)}{d\mu(x)\,d\nu(y)}\Big],
\]
$\Gamma(P,Q)$ being the couplings of $P$ and $Q$.  The primal problem is
\[
  \tag{Primal}
  V \;=\; \sup_{P\in\Pcal(\Zcal)}\Big\{ \E_{z\sim P}[f(z)] \;:\; \W_\epsilon(\widehat\Prob,P)\le\rho \Big\}.
\]

\paragraph{Assumption 1 (verbatim).}
\begin{enumerate}[label=(\Roman*)]
\item $c(x,z)$ is $\widehat\Prob\otimes\nu$-measurable and $\nu\{z: 0\le c(x,z)<\infty\}=1$ for
      $\widehat\Prob$-a.e.\ $x$;
\item $\E_{z\sim\nu}\big[e^{-c(x,z)/\epsilon}\big]<\infty$ for $\widehat\Prob$-a.e.\ $x$;
\item $\Zcal$ is a measurable space and $f:\Zcal\to\R\cup\{\infty\}$ is measurable;
\item every joint $\gamma$ on $\Zcal\times\Zcal$ with first marginal $\widehat\Prob$ admits a regular
      conditional distribution $\gamma_x$.
\end{enumerate}

\paragraph{The shifted radius and the Gibbs kernel (paper's eq.~(2), (3)).}
This is the pivot of the whole theorem, and the point our formalization must track:
\begin{equation}\label{eq:rhobar}
  \boxed{\;\overline{\rho} \;:=\; \rho \;+\; \epsilon\,\E_{x\sim\widehat\Prob}
     \Big[\log \E_{z\sim\nu}\big[e^{-c(x,z)/\epsilon}\big]\Big]\;}
  \qquad
  d\Qcal_{x,\epsilon}(z) \;:=\; \frac{e^{-c(x,z)/\epsilon}}{\E_{u\sim\nu}[e^{-c(x,u)/\epsilon}]}\,d\nu(z).
\end{equation}
$\Qcal_{x,\epsilon}$ is the \emph{Gibbs (tilted) conditional}.  Since $c\ge 0$ and $\nu$ is a
probability measure, $\E_\nu[e^{-c/\epsilon}]\le 1$, so the correction term is $\le 0$ and
$\overline\rho\le\rho$.  Write
\[
  F \;:=\; -\,\epsilon\,\E_{x\sim\widehat\Prob}\Big[\log\E_{z\sim\nu}\big[e^{-c(x,z)/\epsilon}\big]\Big]\;\ge\;0,
  \qquad\text{so}\qquad \overline\rho = \rho - F .
\]
$F$ is the \textbf{free energy} of the tilt; it is $0$ only if $c(x,\cdot)=0$ $\nu$-a.e.

\paragraph{Condition 1 (light tail).}
There exists $\lambda>0$ with $\E_{z\sim\Qcal_{x,\epsilon}}\big[e^{f(z)/(\lambda\epsilon)}\big]<\infty$
for $\widehat\Prob$-a.e.\ $x$.

\paragraph{The dual.}  The paper gives two equal forms,
\begin{align}
  V_D &= \inf_{\lambda\ge 0}\;\Big\{\, \lambda\rho \;+\; \lambda\epsilon\,
        \E_{x\sim\widehat\Prob}\Big[\log \E_{z\sim\nu}\big[e^{(f(z)-\lambda c(x,z))/(\lambda\epsilon)}\big]\Big]\Big\},
        \label{eq:dual-nu}\\
  V_D &= \inf_{\lambda\ge 0}\;\Big\{\, \lambda\overline{\rho} \;+\; \lambda\epsilon\,
        \E_{x\sim\widehat\Prob}\Big[\log \E_{z\sim\Qcal_{x,\epsilon}}\big[e^{f(z)/(\lambda\epsilon)}\big]\Big]\Big\},
        \tag{Dual}\label{eq:dual-Q}
\end{align}
identical for \emph{every} $\lambda>0$: in the exponent the multiplier cancels against the cost,
$\frac{f-\lambda c}{\lambda\epsilon}=\frac{f}{\lambda\epsilon}-\frac{c}{\epsilon}$, whence
\[
  \E_\nu\big[e^{(f-\lambda c)/(\lambda\epsilon)}\big]
  \;=\; \E_\nu\big[e^{-c/\epsilon}\big]\cdot\E_{z\sim\Qcal_{x,\epsilon}}\big[e^{f/(\lambda\epsilon)}\big],
\]
so $\lambda\epsilon\log\E_\nu[\cdot] = \lambda\epsilon\log\E_\nu[e^{-c/\epsilon}] + \lambda\epsilon\log\E_{\Qcal_{x,\epsilon}}[e^{f/(\lambda\epsilon)}]$
and the first term is precisely $\lambda(\overline\rho-\rho)$ after taking $\E_{\widehat\Prob}$.  That
is exactly why the paper can state \eqref{eq:dual-Q} with $\overline\rho$ and \eqref{eq:dual-nu} with
$\rho$.  At $\lambda=0$ the dual objective is \emph{defined} as the $\lambda\downarrow0$ limit, which
equals $\operatorname{ess\,sup}_\nu f$; since it is a limit of the $\lambda>0$ values,
$\inf_{\lambda\ge0}=\inf_{\lambda>0}$, so our Lean dual (an infimum over $\lambda>0$) is the same
number.

\section{Theorem 1 (Strong Duality), verbatim}

\begin{quote}
Let $\widehat\Prob\in\Pcal(\Zcal)$ and assume Assumption 1 holds.  Then:
\begin{enumerate}[label=(\Roman*)]
\item (Primal) is feasible if and only if $\overline\rho \ge 0$;
\item whenever $\overline\rho\ge 0$, $V = V_D$;
\item if in addition Condition 1 holds and $\overline\rho>0$, then $V=V_D<\infty$; otherwise
      $V=V_D=\infty$;
\item assume Condition 1 and $\overline\rho>0$.  With $A:=\{z: f(z)=\operatorname{ess\,sup}_\nu f\}$,
      the dual minimizer $\lambda^*=0$ iff $\operatorname{ess\,sup}_\nu f<\infty$ and
      $\overline\rho \ge \epsilon\,\E_{x\sim\widehat\Prob}\big[\log(1/\Prob_{z\sim\Qcal_{x,\epsilon}}\{A\})\big]$.
\end{enumerate}
If $\overline\rho<0$ then by convention $V=-\infty=V_D$.  Hence $V=V_D$ under Assumption 1 alone.
\end{quote}

\noindent\textbf{Every clause is stated in $\overline\rho$, not $\rho$.}  (Verified by reading p.~7 of
the PDF; \texttt{pdftotext} silently drops the macron, which is how an earlier reading of this repo's
came to record the false statement ``feasible iff $\rho\ge 0$''.)

\section{The paper's proof}

\paragraph{Lemma 1 (weak duality).}  Under Assumption 1,
$V \le \inf_{\lambda\ge0}\{\lambda\rho + \lambda\epsilon\,\E_{\widehat\Prob}[\log\E_\nu e^{(f-\lambda c)/(\lambda\epsilon)}]\} = V_D$.
\begin{distillation}
\item Using Definition 1 and Assumption 1(IV), disintegrate $\gamma$ as
      $d\gamma(x,z)=d\widehat\Prob(x)\,d\gamma_x(z)$ and rewrite
      \[
        V=\sup_{\{\gamma_x\}\subset\Pcal(\Zcal)}\Big\{\E_{\widehat\Prob}\E_{\gamma_x}[f]
        \;:\; \E_{\widehat\Prob}\E_{\gamma_x}\big[c(x,z)+\epsilon\log\tfrac{d\gamma_x(z)}{d\nu(z)}\big]\le\rho\Big\}.
      \]
\item Introduce the multiplier $\lambda\ge0$ for the single scalar constraint:
      $V=\sup_{\{\gamma_x\}}\inf_{\lambda\ge0}\big\{\lambda\rho+\E_{\widehat\Prob}\E_{\gamma_x}
      [f-\lambda c-\lambda\epsilon\log\frac{d\gamma_x}{d\nu}]\big\}$ (paper's eq.~(6)).
\item \emph{Swap} $\sup$ and $\inf$ (this is the inequality --- max-min $\le$ min-max), giving eq.~(7).
\item Push the $\sup$ inside the expectation over $x$ and define the per-point value
      \[
        v_x(\lambda) \;:=\; \sup_{\gamma_x\in\Pcal(\Zcal)}
        \E_{z\sim\gamma_x}\Big[f(z)-\lambda c(x,z)-\lambda\epsilon\log\tfrac{d\gamma_x(z)}{d\nu(z)}\Big].
        \tag{8}
      \]
      Measurability of $x\mapsto v_x(\lambda)$ is Lemma EC.3.
\item \textbf{Gibbs variational principle} (Lemma EC.2): if Condition 1 holds at some $\lambda>0$,
      \[
        v_x(\lambda)\;=\;\lambda\epsilon\,\log\E_{z\sim\nu}\big[e^{(f(z)-\lambda c(x,z))/(\lambda\epsilon)}\big]\;<\;\infty ,
      \]
      the supremum being attained at the tilted measure
      $d\gamma_x^*\propto e^{(f-\lambda c)/(\lambda\epsilon)}d\nu$.  If Condition 1 fails, both sides are
      $+\infty$ and weak duality is vacuous but true.
\end{distillation}

\paragraph{Lemma 2 (reformulation of the primal).}  Under Assumption 1,
\[
  V=\sup_{\{\gamma_x\}\subset\Pcal(\Zcal)}\Big\{\E_{\widehat\Prob}\E_{\gamma_x}[f]
  \;:\; \epsilon\,\E_{x\sim\widehat\Prob}\Big[\E_{z\sim\gamma_x}\log\tfrac{d\gamma_x(z)}{d\Qcal_{x,\epsilon}(z)}\Big]\le\overline\rho\Big\}
  \;=\;\sup\Big\{\cdots:\ \epsilon\,\E_{\widehat\Prob}\big[\KL(\gamma_x\,\|\,\Qcal_{x,\epsilon})\big]\le\overline\rho\Big\}.
\]
\emph{This is the whole content of the $\overline\rho$ shift.}  Substituting \eqref{eq:rhobar} into
$\KL(\gamma_x\|\Qcal_{x,\epsilon}) = \KL(\gamma_x\|\nu) + \tfrac1\epsilon\E_{\gamma_x}[c(x,\cdot)] + \log\E_\nu e^{-c(x,\cdot)/\epsilon}$
turns Definition 1's constraint into the displayed one.  Consequently:

\begin{center}\fbox{\parbox{0.92\textwidth}{
\textbf{Feasibility (Theorem 1(I)) is exactly non-negativity of a relative entropy.}
The constraint functional $\gamma\mapsto \epsilon\,\E_{\widehat\Prob}[\KL(\gamma_x\|\Qcal_{x,\epsilon})]$
is $\ge 0$, with equality \emph{iff} $\gamma_x=\Qcal_{x,\epsilon}$ $\widehat\Prob$-a.s.  So the primal is
feasible iff $\overline\rho\ge0$, and at $\overline\rho=0$ the ambiguity set is the \emph{singleton}
$\{\Prob_0\}$, $d\Prob_0(z)=\E_{x\sim\widehat\Prob}[d\Qcal_{x,\epsilon}(z)]$.
}}\end{center}

\paragraph{Strong duality and the worst-case law.}
Under Condition 1 and $\overline\rho>0$ the paper shows (Lemma 3) that an optimal multiplier
$\lambda^*>0$ exists and is unique, and (Remark 4) the worst-case conditional solves a strictly
convex program with density, w.r.t.\ $\nu$,
\[
  \frac{d\gamma^*_x}{d\nu}(z) \;=\; \alpha_x\,\exp\Big(\big(f(z)-\lambda^*c(x,z)\big)/(\lambda^*\epsilon)\Big),
  \qquad
  \alpha_x=\Big(\E_{z\sim\nu}\big[e^{(f(z)-\lambda^*c(x,z))/(\lambda^*\epsilon)}\big]\Big)^{-1},
\]
so the worst-case law is $\frac{d\Prob_*}{d\nu}(z)=\E_{x\sim\widehat\Prob}\big[\alpha_x e^{(f(z)-\lambda^*c(x,z))/(\lambda^*\epsilon)}\big]$
--- supported on all of $\operatorname{supp}\nu$, unlike Wasserstein DRO's finitely-supported worst case.

\paragraph{The boundary case $\overline\rho=0$.}  Theorem 1(II) asserts $V=V_D$ there too.  It cannot
come from a multiplier: at $\overline\rho=0$ the feasible set is the singleton $\{\Prob_0\}$, so
$V=\E_{\widehat\Prob}\E_{\Qcal_{x,\epsilon}}[f]$, while
$V_D=\inf_{\lambda\ge0}\lambda\epsilon\,\E_{\widehat\Prob}[\log\E_{\Qcal_{x,\epsilon}}e^{f/(\lambda\epsilon)}]$.
By Jensen the dual objective is $\ge \E_{\widehat\Prob}\E_{\Qcal_{x,\epsilon}}[f]$ for every $\lambda>0$,
and $\lambda\epsilon\log\E_{\Qcal}[e^{f/(\lambda\epsilon)}]\downarrow \E_{\Qcal}[f]$ as
$\lambda\to\infty$.  So the infimum is attained \emph{in the limit $\lambda\to\infty$}, not at a finite
multiplier.  \textbf{This limiting argument, not a supergradient, is what carries the boundary.}

\section{Where our Lean development differs}\label{sec:diff}

\begin{formalizationnotes}
\item \textbf{Definition 1 matches.}  \texttt{ForMathlib.OT.Wkappa} $=\W_\epsilon$ with
      $\kappa=\epsilon$, $\mu=\widehat\Prob$, external reference $\nu$; \texttt{sinkhornObjective}
      is the bracket $\E_\gamma[c]+\kappa H(\gamma\mid\widehat\Prob\otimes\nu)$.  Our
      \texttt{sinkhornBall} is Definition 1's ball, i.e.\ stated in the \emph{unbarred} $\rho$.
\item \textbf{The $\overline\rho$ shift is real, and we already found it.}  Our
      \texttt{WangGaoXie2023.primal\_feasible\_radius\_nonneg} records that
      ``$\text{ball nonempty}\iff \rho\ge0$'' is false in Definition-1 coordinates; the true threshold
      is $\rho\ge F$.  Reading p.~7 confirms the paper agrees: its Theorem 1(I) is
      $\overline\rho\ge0$, and $\overline\rho\ge0\iff\rho\ge F$.  \emph{No disagreement} --- the earlier
      repo note calling Theorem 1(I) mis-stated was itself the artifact of a macron-losing text
      extraction.  Corrected here.
\item \textbf{We proved strictly less than the paper at the boundary.}  Our
      \texttt{WangGaoXie2023.sinkhornDual\_le\_droValue} assumes \emph{Slater}: some coupling with
      objective $<\varepsilon$, i.e.\ $\overline\rho>0$.  The paper's Theorem 1(II) claims $V=V_D$ for
      $\overline\rho\ge0$, boundary included.  Our supergradient/concavity assembly genuinely cannot
      reach $\overline\rho=0$ (there is no interior point), and the paper does not use a supergradient
      there --- it uses the $\lambda\to\infty$ Jensen limit above.  \textbf{This is an open item, not a
      refutation:} the boundary case is provable by the limiting argument, and doing so would delete
      the Slater hypothesis from \texttt{strong\_duality} and \texttt{sdrsb\_strong\_duality}.
\item \textbf{Our Slater receipt is not tight.}
      \texttt{ForMathlib.OT.exists\_slater\_of\_couplingCost\_prodCoupling\_lt} discharges Slater from
      $\varepsilon>\E_{\widehat\Prob\otimes\nu}[c]$, using the \emph{independent} coupling.  The tight
      threshold is the free energy $F\le \E_{\widehat\Prob\otimes\nu}[c]$ (Jensen), attained by the
      \emph{Gibbs} coupling $\gamma_x=\Qcal_{x,\epsilon}$.  Sharpening the receipt to $\varepsilon>F$
      is mechanical --- we already have the tilted kernel
      (\texttt{ForMathlib.MeasureTheory.tiltedKernel}) that realizes it.
\item \textbf{Assumption 1(IV) vs.\ our typeclass.}  The paper assumes regular conditional
      distributions exist; we require \texttt{StandardBorelSpace} (and
      \texttt{CountableOrCountablyGenerated} for the KL chain rule), which is the same content in
      Mathlib's idiom, and is what \texttt{Measure.condKernel} needs.
\item \textbf{Condition 1 vs.\ our \texttt{h\_exp}/\texttt{hI\_lp}.}  Condition 1 is
      $\E_{\Qcal_{x,\epsilon}}[e^{f/(\lambda\epsilon)}]<\infty$ for \emph{some} $\lambda>0$; our
      \texttt{h\_exp} asks $\E_\nu[e^{(f-\lambda c)/(\lambda\kappa)}]<\infty$ for \emph{every}
      $\lambda>0$.  These are equivalent up to the tilt's normalizer, but ours is quantified over all
      $\lambda>0$ and is therefore \emph{stronger}.  The paper's Theorem 1(III) is what needs the
      one-$\lambda$ form; our duality proof evaluates the dual at $\lambda^*+\eta$ and so wants the
      family.  Not a fidelity bug, but a hypothesis we could weaken.
\item \textbf{The dual we state is \eqref{eq:dual-nu}, the $\nu$/unbarred-$\rho$ form}
      (\texttt{sinkhornDualObjective} $=\lambda\varepsilon+\E_{\widehat\Prob}[\texttt{logPartition}]$),
      which is the paper's own eq.~(1).  We never form $\Qcal_{x,\epsilon}$ in the dual, so the
      $\overline\rho$ bookkeeping never enters our statement --- it enters only through Slater.
\end{formalizationnotes}

\TODO{Read supplement Lemmas EC.2, EC.3, EC.4 and the proof of Lemma 3 (existence/uniqueness of
$\lambda^*$), and record whether EC.4's $\overline\rho<0\Rightarrow V=V_D=-\infty$ convention is
compatible with our \texttt{sInf}-over-$\lambda>0$ dual (which is a junk value on an empty/unbounded-below
set --- we derive \texttt{BddBelow} from weak duality instead).}

\begin{thebibliography}{99}
\bibitem{ChenGeorgiouPavon2021} Yongxin Chen, Tryphon T. Georgiou, and Michele Pavon. \emph{Stochastic Control Liaisons: Richard Sinkhorn Meets Gaspard Monge on a Schr\"odinger Bridge}. SIAM Review, 63(2):249--313, 2021. doi:10.1137/20M1339982. arXiv:2005.10963.
\bibitem{Leonard2014Survey} Christian L\'eonard. \emph{A Survey of the Schr\"odinger Problem and Some of Its Connections with Optimal Transport}. Discrete and Continuous Dynamical Systems--A, 34(4):1533--1574, 2014. doi:10.3934/dcds.2014.34.1533. arXiv:1308.0215.
\bibitem{Fortet1940} Robert Fortet. \emph{Resolution d'un systeme d'equations de M. Schrodinger}. Journal de mathematiques pures et appliquees, 9e serie, 19(1--4):83--105, 1940.
\bibitem{SinkhornKnopp1967} Richard Sinkhorn and Paul Knopp. \emph{Concerning Nonnegative Matrices and Doubly Stochastic Matrices}. Pacific Journal of Mathematics, 21(2):343--348, 1967.
\bibitem{FranklinLorenz1989} Joel Franklin and Jens Lorenz. \emph{On the Scaling of Multidimensional Matrices}. Linear Algebra and its Applications, 114/115:717--735, 1989. doi:10.1016/0024-3795(89)90490-4.
\bibitem{Carroll2004} Joseph E. Carroll. \emph{Birkhoff's Contraction Coefficient}. Linear Algebra and its Applications, 389:227--234, 2004. doi:10.1016/j.laa.2004.02.039.
\bibitem{EvesonNussbaum1995} Simon P. Eveson and Roger D. Nussbaum. \emph{An Elementary Proof of the Birkhoff--Hopf Theorem}. Mathematical Proceedings of the Cambridge Philosophical Society, 117(1):31--55, 1995.
\bibitem{BlanchetMurthy2019} Jose Blanchet and Karthyek R. A. Murthy. \emph{Quantifying Distributional Model Risk via Optimal Transport}. Mathematics of Operations Research, 44(2):565--600, 2019. doi:10.1287/moor.2018.0936. arXiv:1604.01446.
\bibitem{GaoKleywegt2023} Rui Gao and Anton J. Kleywegt. \emph{Distributionally Robust Stochastic Optimization with Wasserstein Distance}. Mathematics of Operations Research, 48(2):603--655, 2023. doi:10.1287/moor.2022.1275. arXiv:1604.02199.
\bibitem{MohajerinEsfahaniKuhn2018} Peyman Mohajerin Esfahani and Daniel Kuhn. \emph{Data-Driven Distributionally Robust Optimization Using the Wasserstein Metric: Performance Guarantees and Tractable Reformulations}. Mathematical Programming, 171(1--2):115--166, 2018. doi:10.1007/s10107-017-1172-1. arXiv:1505.05116.
\bibitem{WangGaoXie2025} Jie Wang, Rui Gao, and Yao Xie. \emph{Sinkhorn Distributionally Robust Optimization}. Operations Research, 2025. doi:10.1287/opre.2023.0294. arXiv:2109.11926.
\bibitem{Girsanov1960} I. V. Girsanov. \emph{On Transforming a Certain Class of Stochastic Processes by Absolutely Continuous Substitution of Measures}. Theory of Probability and Its Applications, 5(3):285--301, 1960. doi:10.1137/1105027.
\bibitem{CameronMartin1945} R. H. Cameron and W. T. Martin. \emph{Transformations of Wiener Integrals under a General Class of Linear Transformations}. Transactions of the American Mathematical Society, 58(2):184--219, 1945.
\bibitem{Yeh1978} J. Yeh. \emph{Transformation of Conditional Wiener Integrals under Translation and the Cameron--Martin Translation Theorem}. Tohoku Mathematical Journal, 30(4):505--515, 1978.
\bibitem{Kakutani1948} Shizuo Kakutani. \emph{On Equivalence of Infinite Product Measures}. Annals of Mathematics, 49(1):214--224, 1948.
\bibitem{ShiDeBortoliCampbellDoucet2023} Yuyang Shi, Valentin De Bortoli, Andrew Campbell, and Arnaud Doucet. \emph{Diffusion Schr\"odinger Bridge Matching}. Advances in Neural Information Processing Systems, 2023. arXiv:2303.16852.
\bibitem{LiuLipmanNickelKarrerTheodorouChen2024} Guan-Horng Liu, Yaron Lipman, Maximilian Nickel, Brian Karrer, Evangelos A. Theodorou, and Ricky T. Q. Chen. \emph{Generalized Schr\"odinger Bridge Matching}. International Conference on Learning Representations, 2024. arXiv:2310.02233.
\bibitem{DomingoEnrichDrozdzalKarrerChen2025} Carles Domingo-Enrich, Michal Drozdzal, Brian Karrer, and Ricky T. Q. Chen. \emph{Adjoint Matching: Fine-Tuning Flow and Diffusion Generative Models with Memoryless Stochastic Optimal Control}. arXiv preprint arXiv:2409.08861, 2025.

\bibitem{Birkhoff1957} Garrett Birkhoff. \emph{Extensions of Jentzsch's Theorem}. Transactions of the American Mathematical Society, 85(1):219--227, 1957.
\bibitem{Sinkhorn1964} Richard Sinkhorn. \emph{A Relationship Between Arbitrary Positive Matrices and Doubly Stochastic Matrices}. Annals of Mathematical Statistics, 35(2):876--879, 1964.
\bibitem{Sinkhorn1967Prescribed} Richard Sinkhorn. \emph{Diagonal Equivalence to Matrices with Prescribed Row and Column Sums}. American Mathematical Monthly, 74(4):402--405, 1967.
\bibitem{Bauer1965} Friedrich L. Bauer. \emph{An Elementary Proof of the Hopf Inequality for Positive Operators}. Numerische Mathematik, 7:331--337, 1965.
\bibitem{Bushell1973} P. J. Bushell. \emph{Hilbert's Metric and Positive Contraction Mappings in a Banach Space}. Archive for Rational Mechanics and Analysis, 52:330--338, 1973.

\end{thebibliography}

\end{document}
